% !Tex root = ../main.tex

Die Abbildung~\ref{fig:cc-rateExec} zeigt den Einfluss von Concurrency auf die Erfolgs-/Fehlerraten und auf die
Ausführungszeiten der Transaktionen.
Über alle Fehlerklassen sinkt die Erfolgsrate bei höherer Konkurrenz.
Bei Lock Timeout Fehlern ist die Erfolgschance über alle Concurrency Werte die Geringste.
Sie startet bei 36\% und fällt exponentiell auf 3\% bei Concurrency von 20.
Die höchste Erfolgschance haben Serialization Failures die bei 95\% starten und bei 26\% enden.
Den stärksten Abfall haben Deadlocks, die auch bei 95\% starten aber bis auf 10\% sinken.
Die Ausführungszeiten von Serialization Failures sind die geringsten.
Dabei ist es interessant zu sehen, dass bis zu einer Concurrency von 7 die Ausführungszeit von erfolgreichen
Transaktionen über der von nicht erfolgreichen Transaktionen liegt.
Ein Unterschied wird erst ab 10 gleichzeitigen Anwendungen erkennbar.
Des Weiteren lässt sich festellen, dass die Ausführungszeit von erfolgreichen Transaktionen bei höherer Concurrence
wieder sinkt, während sie bei allen anderen stagniert oder steigt.
Die Ausführungszeiten bei Lock Timeouts laufen parallel zueinander, wobei der Abstand zwischen diesen darauf
zurückzuführen ist, dass eine Transaktion wegen der Implementierung mindestens eine Sekunde warten muss.
In der Zeit kann eine weitere schon abgebrochen werden.
Die höchste Ausführungszeit haben Deadlocks.
Bei einer Concurrency von 2 liegt diese noch unter 2000ms und steigt dann stark an auf über 15000ms.
Erfolgreiche und nicht erfolgreiche Transaktionen gleichen sich bei steigender Concurrency an und bei Concurrency
von 20 haben erfolgreiche Transaktion im Durchschnitt länger gebraucht.
Die durchschnittliche Dauer einer erfolgreichen Transaktion liegt bei 16841ms, während die von nicht erfolgreichen
Transaktionen bei 15450ms liegt.

\input{figures/cc-rateExec}

Interessant ist außerdem, dass sowohl bei Deadlocks, als auch bei Serialization Failures für eine Concurrency von 2
nicht mehr als 3 Retries benötigt wurden, wie die Tabelle~\ref{tab:cc-retryDistribution} zeigt.
Die in Tabelle~\ref{tab:cc-quantiles} gezeigten Quantile zeigen auch eine Besonderheit.
Während bei Deadlocks und Lock Timeouts die Ausführungszeit über die Quantile und über die Concurrency steigt, hat bei
Serialization die Concurrency von 5 die längste Ausführungszeit in dem Quantil $p_{99.9}$.
Dies ist bei erfolgreichen und nicht erfolgreichen Transaktionen der Fall.
In diesem Quantil steigt die Ausführungszeit erst mit höheren Concurrency Werten und dann sinkt sie wieder ab
einer Concurrency von 10.

Der Retry Overhead~\ref{fig:cc-overhead} verhält sich analog zur Ausführungszeit, außer das die Werte niedriger sind
und bei Concurrency von 2 der durchschnittliche Overhead für Deadlocks größer ist als der von Lock Timeouts.
Bei der Ausführungszeit war dies noch andersherum.
Der Durchsatz sinkt mit höherer Concurrency wie Abbildung~\ref{fig:cc-throughput} zeigt.
Dabei ist interessant, dass bei 20 Concurrency der Durchsatz ohne Retry höher ist als mit Retry.
Dies ist aber nur bei Serialization Failure und Lock Timeout der Fall.
Bei Deadlocks hingegen hat \texttt{Experiment0} den höchsten Durchsatz bei Concurrency von 2 und den schlechtesten 20.
Generell ist der Durchsatz von Deadlocks bei jedem Concurrency Level niedriger als der, der anderen beiden
Fehlerklassen.
Serialization Failures haben den höchsten Durchsatz.